\section{Model Training and Dataset}\label{sec:training}

Training a reliable aerial object detector for SAR requires data that closely represents the deployment domain---a mannequin viewed from altitude against a grassy field background. Obtaining large volumes of real labelled aerial imagery is expensive and weather-dependent, so a mixed-source training strategy was adopted that combines synthetic composites, a small set of real labelled frames, and hard negative images. This section describes the dataset construction pipeline, training configuration, and deployment workflow.

\subsection{Training Data Strategy}\label{sec:training:strategy}

The final dataset (\texttt{dataset\_v2}) contains 366 images at the native camera resolution of $1456\times1088$ pixels drawn from three complementary sources:

\begin{enumerate}
    \item \textbf{Synthetic composites} (300 images). A cutout of the mannequin (\texttt{dummy.png}, with alpha channel) is composited onto random crops from real DJI flight video backgrounds at altitude-correct scales. These images provide the bulk of the training volume and allow systematic variation of target size, orientation, and position.
    \item \textbf{Real labelled frames} (16 images). Frames extracted from a DJI flight recording at multiple altitudes (15--50\,m) and manually annotated with bounding boxes. These anchor the model to real-world textures, lighting, and target appearance.
    \item \textbf{Negative images} (50 images). Background frames from the same flight video containing no mannequin, paired with empty label files. These teach the model to suppress false positives on terrain features such as shadows, tarpaulins, and path markings.
\end{enumerate}

This three-source approach follows the domain-randomisation principle~\cite{domainrand}: synthetic data provides scale and diversity, real data provides domain fidelity, and negatives provide specificity. Training at the native sensor resolution rather than the standard $640\times640$ YOLO input size allows the network to learn fine-grained features during training, even though inference-time resizing reduces the input to $640\times640$~\cite{yolov8}.

\subsection{Synthetic Data Generation}\label{sec:training:synthetic}

Synthetic images are generated by \texttt{generate\_dataset\_v2.py}, which performs the following steps for each of the 300 output images:

\begin{enumerate}
    \item A random $1456\times1088$ crop is extracted from a high-resolution satellite image (\texttt{map.jpg}) or from DJI video frames, providing diverse background textures.
    \item The mannequin cutout is scaled to simulate a target altitude drawn from a three-tier distribution: 30\% close (scale 0.20--0.40, simulating 15--20\,m), 40\% medium (0.08--0.20, simulating 20--35\,m), and 30\% far (0.03--0.08, simulating 35--60\,m). This biases training towards the operationally critical medium-altitude band while still covering edge cases.
    \item The scaled dummy is rotated by a random angle ($0$--$360^{\circ}$) and alpha-blended onto a random position within the background crop.
    \item Data augmentation is applied stochastically: brightness and contrast jitter (50\% probability, $\alpha \in [0.7, 1.3]$, $\beta \in [-30, 30]$) and Gaussian blur to simulate motion blur (30\% probability, kernel $3\times3$ or $5\times5$).
    \item A YOLO-format label file is written containing the normalised bounding box coordinates (class, $x_{\text{centre}}$, $y_{\text{centre}}$, width, height).
\end{enumerate}

Compared to the v1 dataset (200 images at $640\times640$ from a single satellite background), v2 introduces real-flight backgrounds, altitude-correct target scaling, a negative set, and native-resolution output---all of which reduce the domain gap between training and deployment~\cite{synthdata}.

\subsection{Real Data Labelling}\label{sec:training:labelling}

A custom labelling tool (\texttt{label\_tool.py}) was developed for rapid bounding-box annotation of video frames. The tool displays each frame at native resolution and supports click-to-place annotation: left-click sets the box centre, scroll wheel adjusts box size, \texttt{S} saves, and \texttt{N} skips. The \texttt{-{}-full} flag ensures labels are computed relative to the full $1456\times1088$ frame rather than $640\times640$ tiles, preserving spatial fidelity.

Sixteen frames were labelled from the DJI flight recording \texttt{DJI\_0001\_1456x1088\_cropped\_30fps.mp4}, selected to span a range of altitudes (15--50\,m) and viewing angles. While small in number, these real frames proved disproportionately valuable: they contain authentic lighting, lens effects, and background clutter that synthetic composites cannot fully replicate.

\subsection{Negative Mining}\label{sec:training:negatives}

Fifty frames were extracted from the same flight video at locations where no mannequin was present. Each negative image is paired with an empty label file, signalling to the YOLO loss function that no objects exist in the frame. Including negatives is a standard technique for reducing false-positive rates in object detection~\cite{hardneg}, and proved essential for this project: the v1 model (trained without negatives) frequently hallucinated detections on shadows and field markings, whereas the v2 model suppresses these spurious activations.

\subsection{Training Pipeline}\label{sec:training:pipeline}

Training was performed on Google Colab using an NVIDIA T4 GPU. The model was initialised from COCO-pretrained YOLOv8n weights (\texttt{yolov8n.pt}), applying transfer learning to leverage low-level feature representations learned from the 80-class COCO dataset~\cite{coco_dataset, yolov8}. Table~\ref{tab:training_config} summarises the training configuration.

\begin{table}[ht]
    \centering
    \caption{Training hyperparameters for the YOLOv8n SAR dummy detector (v2).}
    \label{tab:training_config}
    \begin{tabular}{ll}
        \toprule
        \textbf{Parameter} & \textbf{Value} \\
        \midrule
        Base model & YOLOv8n (COCO-pretrained) \\
        Training image size & $1088 \times 1088$ \\
        Epochs & 150 (early stopping patience: 30) \\
        Batch size & 8 \\
        Augmentation: rotation & $\pm20^{\circ}$ \\
        Augmentation: scale & 0.3--1.7 (aggressive, simulating altitude) \\
        Augmentation: mosaic & 1.0 (enabled throughout) \\
        Augmentation: mixup & 0.1 \\
        Augmentation: copy-paste & 0.2 \\
        Horizontal flip & 0.5 \\
        \bottomrule
    \end{tabular}
\end{table}

The training-time augmentation parameters were tuned for aerial detection: aggressive scale variation ($0.7$) simulates altitude changes, mosaic augmentation improves small-object detection by tiling four images into one training sample~\cite{yolov4}, and copy-paste augmentation randomly pastes annotated objects onto new backgrounds for additional diversity. Random erasing was explicitly disabled (\texttt{erasing=0.0}) because it risks erasing the small target entirely.

\subsection{Results}\label{sec:training:results}

Table~\ref{tab:training_results} presents the validation metrics for the v2 model compared to the original v1 baseline. The v2 model achieves near-perfect detection performance on the validation set.

\begin{table}[ht]
    \centering
    \caption{Detection performance comparison between model versions. All metrics are computed on the respective validation sets.}
    \label{tab:training_results}
    \begin{tabular}{lcccc}
        \toprule
        \textbf{Model} & \textbf{mAP\textsubscript{50}} & \textbf{mAP\textsubscript{50--95}} & \textbf{Precision} & \textbf{Recall} \\
        \midrule
        v1 (200 syn, 640) & $\sim$0.95 & --- & --- & --- \\
        v2 (366 mixed, 1088) & \textbf{0.995} & 0.828 & 0.994 & 0.989 \\
        \bottomrule
    \end{tabular}
\end{table}

A caveat applies: the training and validation splits use the same image set (\texttt{train=val=images}), which inflates the reported metrics relative to a held-out evaluation. This pragmatic choice was driven by the small dataset size---splitting 366 images 80/20 would leave only 73 validation images, too few for stable metric estimation. The strong field-test results (Section~\ref{sec:field}) provide independent confirmation that the model generalises beyond the training distribution.

The exported TFLite model is 11.7\,MB (float32) and maintains the same input tensor shape $[1, 640, 640, 3]$ and output tensor shape $[1, 5, 8400]$ as the original model. Three model variants are archived in the \texttt{cv\_models/} directory: \texttt{sar\_v2\_1088} (current best), \texttt{sar\_640}, and \texttt{sar\_1280}, each containing TFLite, PyTorch, and NCNN exports along with training curves and confusion matrices for reproducibility.

\subsection{Deployment}\label{sec:training:deployment}

The deployment workflow is deliberately minimal. Because all model variants share the same input/output tensor shapes, swapping the active model on the Raspberry Pi requires a single file copy:

\begin{verbatim}
cp cv_models/sar_v2_1088/best.tflite best.tflite
\end{verbatim}

No code changes, configuration edits, or dependency updates are needed. The \texttt{vision.py} module loads whichever \texttt{best.tflite} file is present in the project root, and all downstream modules consume the same \texttt{(found, cx, cy, confidence)} interface regardless of which model generated the detection. This design supports rapid model iteration on flight day: if a model produces excessive false positives in the field, an alternative can be deployed in under 30 seconds by copying a different TFLite file from the \texttt{cv\_models/} directory.

An NCNN export is also produced alongside each TFLite model, offering a potential path to approximately $3\times$ faster inference on the Pi~5 ARM cores if future mission requirements demand higher frame rates. The \texttt{-{}-model} flag in \texttt{main.py} allows runtime model selection without file copying, further streamlining field testing of alternative models.

\subsection{Limitations}\label{sec:training:limitations}

Several limitations of the current training pipeline should be acknowledged:

\begin{itemize}
    \item \textbf{Train/validation overlap.} As noted in Section~\ref{sec:training:results}, the \texttt{dataset.yaml} configuration uses the same image directory for both training and validation. The reported mAP$_{50}$ of 0.995 therefore reflects in-sample performance and likely overstates true generalisation accuracy. A proper stratified 80/20 split would provide more meaningful validation metrics, but was not pursued due to the small dataset size.

    \item \textbf{Small real-image fraction.} Only 16 of 366 images (4.4\%) are real aerial photographs; the remaining 95.6\% are synthetic composites or negatives derived from a single satellite image and one flight video. While domain randomisation mitigates this to some extent~\cite{domainrand}, the model may struggle with target appearances, lighting conditions, or ground textures not represented in the limited real data.

    \item \textbf{Single target appearance.} All synthetic images use the same mannequin foreground image. A real SAR scenario would present casualties in varied clothing, body positions, and degrees of occlusion. The model's ability to generalise to human targets beyond the specific mannequin used in training has not been validated.

    \item \textbf{Limited environmental diversity.} Training backgrounds are drawn from one satellite image and one flight video recorded on a single day. Seasonal variation, different terrain types, adverse weather, and varying solar angles are not represented. Deployment in conditions significantly different from the training environment may degrade detection performance.

    \item \textbf{No partial occlusion.} The dataset contains no examples of a partially occluded target---for instance, a casualty partially hidden by vegetation, rocks, or terrain features---which is a common scenario in real wilderness SAR operations.
\end{itemize}

Despite these limitations, the trained model performed reliably during bench testing and video replay analysis across the altitude range of 15--50\,m (Section~\ref{sec:results}), and the modular deployment architecture allows rapid model replacement as improved datasets become available.
